\documentclass[11pt]{article}
\usepackage{amsmath,amssymb}
\title{Restricted Affine/Horn Elimination on Exact-(2,2)-CNF-3}
\author{Grok 4.7}
\date{31 October 2026, 18:15 PDT}
\begin{document}
\maketitle

\begin{abstract}
We record a restricted exact elimination procedure for pure affine and Horn fragments of exact-(2,2)-CNF-3. The procedure is linear-time and produces an equivalent residual formula. On mixed nonlinear instances the residual core remains arbitrary NP-complete. This restates a known limitation and does not yield a universal resource contraction.
\end{abstract}

\section{Definitions}
An \emph{exact-(2,2)-CNF-3} formula has every clause of width exactly 3 with distinct variables, and every variable occurs exactly twice positively and twice negatively. We consider the fragment in which every clause is either a Horn clause or a parity (XOR) clause.

\section{Restricted Lemma (PROVED/RESTRICTED)}
\begin{theorem}[Degree-2 resolution + Gaussian elimination]
Let $F$ be an exact-(2,2)-CNF-3 formula that is a conjunction of Horn clauses and XOR clauses. There exists a deterministic algorithm that, in linear time and linear space, produces an equivalent residual formula $Q$ on a (possibly empty) subset of the variables such that:
\begin{itemize}
\item every variable of degree 2 is eliminated by resolution or Gaussian elimination;
\item $|Q| = O(|F|)$;
\item $F$ is satisfiable if and only if $Q$ is satisfiable;
\item a satisfying assignment for $Q$ extends to one for $F$ in linear time.
\end{itemize}
\end{theorem}

\begin{proof}[Sketch]
Variables of occurrence degree 2 may be resolved out in linear time by the standard resolution rule for 3-clauses, producing at most a constant number of new clauses per elimination (bounded by the exact-(2,2) condition). Connected XOR components are solved by Gaussian elimination over $\mathrm{GF}(2)$ in linear time (sparse implementation). The process terminates with a residual formula free of degree-2 variables. Correctness follows from the soundness and completeness of resolution and Gaussian elimination on the respective fragments. Size is preserved up to a constant factor because each elimination removes one variable and adds $O(1)$ clauses.
\end{proof}

\section{Limitation (KNOWN)}
On a general mixed formula containing non-Horn, non-XOR clauses (or after the pure fragments are eliminated), the residual core $Q$ can be an arbitrary NP-complete SAT instance of linear size. Expander Tseitin formulas provide concrete examples in which the affine shell is eliminated yet the core remains hard. Consequently the procedure does not reduce the asymptotic size below $N^{1-o(1)}$ on worst-case instances, and therefore cannot by itself force a simultaneous time-space contradiction against the FLMV region for an arbitrary polynomial-time solver $A$.

\section{Status}
The lemma is \textbf{PROVED} for the pure Horn/XOR fragment and \textbf{RESTRICTED} to that fragment. The universal implication from an arbitrary polynomial-time SAT solver to an impossible resource profile remains \textbf{OPEN}.

\end{document}
